\begin{tikzpicture}[line width=0.5pt,font=\footnotesize]
% 第一栏
\draw (0.85,3.42) -- (0.30,2.58);
\draw (0.80,2.62) -- (0.30,1.82);
\draw (0.80,1.86) -- (0.30,1.06);
\draw (0.85,1.10) -- (0.30,0.30);
\node[anchor=west] at (1.00,3.02) {$E\ (-1-c)$};
\node[anchor=west] at (1.00,2.27) {$E_3\ (-3c)$};
\node[anchor=west] at (1.00,1.52) {$E_2\ (-2c)$};
\node[anchor=west] at (1.00,0.77) {$E_1\ (-c)$};
% 第二栏
\draw (3.65,3.42) -- (3.10,2.58);
\draw (3.60,2.62) -- (3.10,1.82);
\draw (3.60,1.86) -- (3.10,1.06);
\node[anchor=west] at (3.80,3.02) {$E\ (-1-c)$};
\node[anchor=west] at (3.80,2.27) {$E_2\ (-2c)$};
\node[anchor=west] at (3.80,1.52) {$E_1\ (-c)$};
% 第三栏
\draw (6.45,3.42) -- (5.90,2.58);
\draw (6.40,2.62) -- (5.90,1.82);
\node[anchor=west] at (6.60,3.02) {$E\ (-1-c)$};
\node[anchor=west] at (6.60,2.27) {$E_1\ (-c)$};
% 第四栏
\draw (9.25,3.42) -- (8.70,0.30);
\node[anchor=west] at (9.40,3.02) {$E\ (-1-c)$};
\fill (9.02,1.90) circle (0.05);
\node[anchor=west] at (9.15,1.80) {$Z_0$ 的泛点};
% 收缩箭头
\draw[->] (2.55,2.45) -- (2.95,2.45);
\node[above] at (2.75,2.45) {$g_3$};
\draw[->] (5.35,2.45) -- (5.75,2.45);
\node[above] at (5.55,2.45) {$g_2$};
\draw[->] (8.15,2.45) -- (8.55,2.45);
\node[above] at (8.35,2.45) {$g_1$};
\end{tikzpicture}
